\documentclass[11pt,a4paper]{article}
\usepackage{amsmath,amssymb,amsthm}
\usepackage{geometry}
\geometry{margin=1in}
\usepackage{hyperref}

\title{\bfseries Glossary of Terminology for the Harmonic--Categorical Proof of the Riemann Hypothesis}
\author{}
\date{}

\begin{document}
\maketitle

\section*{Functions and Series}

\begin{description}
\item[Riemann zeta function] \(\zeta(s) = \sum_{n=1}^\infty n^{-s}\) for \(\operatorname{Re}s>1\), meromorphically continued to \(\mathbb{C}\) with a simple pole at \(s=1\). Its non‑trivial zeros are the subject of the Riemann Hypothesis. \cite{HSTpaper,RHproof}
\item[Greedy harmonic expansion] For \(\theta\in[0,\pi]\), \(\displaystyle \frac{\theta}{\pi} = \sum_{n=0}^\infty \frac{\delta_n(\theta)}{n+2}\), where \(\delta_n(\theta)\in\{0,1\}\) are the digits produced by the greedy step‑function algorithm. \cite{HSTpaper}
\item[Greedy harmonic zeta function] \(\displaystyle \Phi(\theta,s) = \sum_{n=0}^\infty \frac{\delta_n(\theta)}{(n+2)^s}\), \(\operatorname{Re}s>1\). Continued meromorphically to \(\mathbb{C}\) with poles only at the non‑trivial zeros of \(\zeta(s)\) (and a simple pole at \(s=1\)). \cite{HSTpaper}
\item[HST Dirichlet series] For \(f\in L^2[0,\pi]\), \(\displaystyle D_f(s) = \sum_{\psi\in\mathcal{B}} \frac{\mathcal{H}f(\psi)}{(|\psi|+2)^s}\), \(\operatorname{Re}s>1\). Encodes the Harmonic Sine Transform coefficients. \cite{HSTpaper}
\item[Entire factor \(P(s)\)] The entire function appearing in the Fredholm determinant identity: \(\det(I-\mathcal{L}_s) = \frac{\zeta(s)}{\zeta(2s)}e^{P(s)}\). It is non‑vanishing in the critical strip. \cite{DeterminantProof,RHproof}
\end{description}

\section*{Operators and Matrices}

\begin{description}
\item[Transfer operator \(\mathcal{L}_s\)] Acting on \(H^2(\mathbb{D})\), for \(\operatorname{Re}s>1\) defined by \[(\mathcal{L}_s f)(z) = \sum_{j=1}^\infty \frac{j+1}{(j+2)^{s+1}}\Bigl[f\!\bigl(\tfrac{j+1}{j+2}z\bigr) + f\!\bigl(\tfrac{j+1}{j+2}z + \tfrac{1}{j+2}\bigr)\Bigr].\] Trace‑class for \(\operatorname{Re}s>1\), Hilbert–Schmidt on the critical line. \cite{HSTpaper,TraceClass}
\item[Truncated transfer operator \(A_k(t)\)] The \(d_k\times d_k\) matrix \(A_k(t) = P_k\mathcal{L}_{1/2+it}P_k|_{\mathcal{H}_k}\), where \(P_k\) projects onto the first \(d_k = k+1\) monomials. \cite{RHproof}
\item[Hamiltonian \(H_k\)] The Hermitian matrix \(H_k = i\frac{d}{dt}\log U_k(t)|_{t=0}\), with \(U_k(t)\) the unitary factor of \(A_k(t)\) from polar decomposition. \cite{RHproof}
\item[Metric operator \(\eta\)] (Not required in final proof) A hypothetical bounded, positive, invertible operator intertwining the transfer operator and its adjoint, used in early attempts. \cite{SpectralProgramme}
\item[Frame operator \(\mathcal{S}\)] \(\mathcal{S}f = \sum_n \langle f,e_n\rangle e_n\) for the digit functions \(e_n\); shown not to be boundedly invertible. \cite{ResearchGuide}
\end{description}

\section*{Spaces and Bases}

\begin{description}
\item[Hardy space \(H^2(\mathbb{D})\)] The Hilbert space of analytic functions \(f(z)=\sum_{n\ge0}a_n z^n\) on the unit disc with \(\sum|a_n|^2<\infty\). \cite{HSTpaper}
\item[Greedy harmonic tree \(\mathcal{T}\)] Binary tree whose nodes are cylinder intervals of the greedy harmonic expansion; the tree is complete and separates points of \([0,\pi]\). \cite{HSTpaper}
\item[Haar basis \(\mathcal{B}\)] Orthonormal basis of \(L^2[0,\pi]\) consisting of the constant function and the Haar wavelets \(\psi_\varepsilon\) attached to split nodes of \(\mathcal{T}\). \cite{HSTpaper}
\item[Harmonic Sine Transform (HST)] The isometric isomorphism \(\mathcal{H}: L^2[0,\pi] \to \ell^2(\mathcal{B})\), \(\mathcal{H}f(\psi) = \langle f,\psi\rangle\). \cite{HSTpaper}
\end{description}

\section*{Determinants and Spectral Quantities}

\begin{description}
\item[Fredholm determinant \(\Delta(t)\)] \(\Delta(t) = \det_{(2)}(I-\mathcal{L}_{1/2+it})\), regularised determinant on the critical line. Zeros are the ordinates \(\gamma\) of non‑trivial zeros. \cite{RHproof,DeterminantProof}
\item[Truncated determinant \(\Delta_k(t)\)] \(\Delta_k(t) = \det(I-A_k(t))\), a polynomial in the entries of \(A_k(t)\). Its zeros exactly equal the eigenvalues of \(H_k\) (Lemma~4.1). \cite{RHproof,SpectralFlowLemma}
\item[Spectral measure of \(H_k\)] \(\mu_k = \sum_{\lambda\in\operatorname{spec}(H_k)}\delta_\lambda\); converges vaguely to the zero‑counting measure \(\mu = \sum_\gamma m_\gamma \delta_\gamma\). \cite{RHproof}
\item[Zero‑counting measure] \(\mu = \sum_{\gamma:\zeta(1/2+i\gamma)=0} m_\gamma \delta_\gamma\), where \(m_\gamma\) is the multiplicity of the zero. \cite{RHproof}
\item[Regularised determinant \(\det\nolimits_{(2)}\)] Defined for Hilbert–Schmidt operators: \(\det\nolimits_{(2)}(I-T) = \det\nolimits_{(1)}((I-T)e^T)\). \cite{TraceClass}
\end{description}

\section*{Key Theorems and Lemmas}

\begin{description}
\item[Haar Basis Theorem (Thm.~2.1)] The set \(\mathcal{B}\) is an orthonormal basis of \(L^2[0,\pi]\). \cite{HSTpaper}
\item[Meromorphic Continuation (Thm.~3.1)] \(\Phi(\theta,s)\) extends meromorphically to \(\mathbb{C}\) with poles only at the non‑trivial zeros of \(\zeta(s)\) (and \(s=1\)). \cite{HSTpaper}
\item[Fredholm Determinant Identity (Thm.~5.1)] \(\det\nolimits_{(2)}(I-\mathcal{L}_{1/2+it}) = \frac{\zeta(1/2+it)}{\zeta(1+2it)} e^{P(1/2+it)}\), with \(P\) entire and zero‑free on the critical strip. \cite{DeterminantProof,RHproof}
\item[Finite‑Dimensional Spectral Correspondence (Lemma~4.1)] The eigenvalues of the Hermitian matrix \(H_k\) are exactly the real zeros of \(\Delta_k(t)\), with multiplicities. \cite{SpectralFlowLemma,RHproof}
\item[Uniform Convergence of Determinants (Prop.~6.1)] \(\Delta_k(t) \to \Delta(t)\) uniformly on compact subsets of \(\mathbb{R}\). \cite{RHproof}
\item[Hurwitz's Theorem] The zeros of a uniformly convergent sequence of analytic functions converge to zeros of the limit function (counting multiplicities). \cite{RHproof}
\item[Riemann Hypothesis (Thm.~7.1)] All non‑trivial zeros of \(\zeta(s)\) lie on the line \(\operatorname{Re}s = 1/2\). \cite{RHproof}
\end{description}

\section*{Constants and Special Values}

\begin{description}
\item[\(\pi\)] The circle constant; appears as the length of the interval \([0,\pi]\) and in harmonic fractions \(\pi/(n+2)\). \cite{HSTpaper}
\item[Harmonic fractions \(\alpha_n = \frac{\pi}{n+2}\)] The fundamental lengths used in the greedy algorithm; also the right‑child lengths on the tree. \cite{HSTpaper}
\item[Dimensions \(d_k = k+1\)] The size of the finite‑dimensional subspaces \(\mathcal{H}_k\). \cite{RHproof}
\end{description}

\section*{Notation Index}

\begin{center}
\begin{tabular}{llp{8cm}}
\(\theta\) & Angle in \([0,\pi]\) & \\
\(\delta_n(\theta)\) & Greedy digit (0 or 1) at step \(n\) & \cite{HSTpaper} \\
\(r_n(\theta)\) & Remainder after \(n\) steps & \cite{HSTpaper} \\
\(\mathcal{T}\) & Greedy harmonic tree & \cite{HSTpaper} \\
\(\varepsilon\) & Node label (binary prefix) & \cite{HSTpaper} \\
\(\psi_\varepsilon\) & Haar wavelet at node \(\varepsilon\) & \cite{HSTpaper} \\
\(\psi_\varnothing\) & Constant wavelet \(1/\sqrt{\pi}\) & \cite{HSTpaper} \\
\(\mathcal{B}\) & Orthonormal Haar basis & \cite{HSTpaper} \\
\(\mathcal{H}\) & Harmonic Sine Transform (HST) & \cite{HSTpaper} \\
\(\mathcal{L}_s\) & Greedy transfer operator & \cite{HSTpaper} \\
\(A_k(t)\) & Truncated transfer operator on the critical line & \cite{RHproof} \\
\(U_k(t)\) & Unitary factor of \(A_k(t)\) & \cite{RHproof} \\
\(H_k\) & Finite‑dimensional Hermitian Hamiltonian & \cite{RHproof} \\
\(\Delta(t)\) & Full Fredholm determinant on critical line & \cite{RHproof} \\
\(\Delta_k(t)\) & Truncated determinant \(\det(I-A_k(t))\) & \cite{RHproof} \\
\(Z\) & Set of ordinates of non‑trivial zeros of \(\zeta(s)\) & \\
\(Z_k\) & Zero set of \(\Delta_k(t)\) (eigenvalues of \(H_k\)) & \cite{RHproof} \\
\(\mu\) & Zero‑counting measure & \cite{RHproof} \\
\(P(s)\) & Entire factor in determinant identity & \cite{DeterminantProof} \\
\(\operatorname{spec}(H)\) & Spectrum of operator \(H\) & \\
\(\det\nolimits_{(1)},\det\nolimits_{(2)}\) & Ordinary and regularised Fredholm determinants & \cite{TraceClass} \\
\end{tabular}
\end{center}

\section*{Conjectures}

\begin{description}
\item[Hilbert–Pólya conjecture] There exists a self‑adjoint operator whose spectrum coincides with the non‑trivial zeros of \(\zeta(s)\). The present work realises this conjecture. \cite{RHproof,SpectralProgramme}
\end{description}

\begin{thebibliography}{99}

\bibitem{HSTpaper}
V.~Geere et al., \emph{A Harmonic Sine Transform and the Riemann Zeta Function}, 2026.

\bibitem{RHproof}
V.~Geere et al., \emph{A Rigorous Proof of the Riemann Hypothesis via the Harmonic–Categorical Hilbert–Pólya Operator}, 2026.

\bibitem{ResearchGuide}
V.~Geere et al., \emph{Foundational Problems for the Hilbert–Pólya Operator}, research guide, 2026.

\bibitem{SpectralProgramme}
V.~Geere et al., \emph{Spectral Analysis of the Limiting Hilbert–Pólya Operator: A Research Programme}, 2026.

\bibitem{DeterminantProof}
V.~Geere et al., \emph{Rigorous Fredholm determinant identity for the greedy harmonic transfer operator}, companion paper, 2026.

\bibitem{TraceClass}
V.~Geere et al., \emph{Trace‑class extension and regularised determinants}, companion paper, 2026.

\bibitem{SpectralFlowLemma}
V.~Geere, \emph{Finite‑dimensional spectral flow and the eigenvalue–determinant correspondence}, supplementary note, 2026.

\end{thebibliography}

\end{document}